\section{সরকারি কাজে Database}

Corporate Database ও Government Database পৃথক Data Model নয়; এগুলো ব্যবহারের ক্ষেত্রভিত্তিক নাম।

\subsubsection{আবেদন জমা দেওয়ার পরে নাগরিক জানবেন কীভাবে?}
একটি সেবাকেন্দ্রে অনেক আবেদন জমা পড়েছে। কোনো আবেদনের কোন নথি বাকি এবং সেটি কোন কর্মীর কাছে আছে, তা জানতে খাতা খুঁজতে হয়। নাগরিক বারবার অগ্রগতি জানতে এলে কর্মীদেরও একই তথ্য বারবার খুঁজতে হয়।

এই সমস্যা কমাতে প্রতিটি আবেদনের ApplicationID, বর্তমান অবস্থা, সময় ও দায়িত্বপ্রাপ্ত শাখার তথ্য সংরক্ষণ করা প্রয়োজন। অনুমোদিত কর্মী এসব তথ্য ব্যবহার করে নির্দিষ্ট আবেদন খুঁজে তার অগ্রগতি জানাতে পারেন। অ্যাপ্লিকেশনের মাধ্যমে নাগরিকও নিজের আবেদনের অবস্থা দেখতে পারেন। নিবন্ধন, সনদ ও শিক্ষাসেবার আবেদনে এই পদ্ধতি উপযোগী। এতে তথ্য খোঁজা সহজ হয়; তবে অসম্পূর্ণ নথি বৈধ হয়ে যায় না।

\begin{definitionbox}{Government Database}
জনসেবা, প্রশাসন, পরিকল্পনা ও সরকারি কার্যক্রমের তথ্য সংরক্ষণ ও ব্যবহারের জন্য সরকারি প্রতিষ্ঠানে ব্যবহৃত ডেটাবেজকে সরকারি ডেটাবেজ বলে।
\end{definitionbox}
নাগরিকের পরিচয়, জন্ম-মৃত্যু নিবন্ধন, ভূমির রেকর্ড, শিক্ষা, স্বাস্থ্য, রাজস্ব ও জনশুমারির তথ্য এর পরিচিত ক্ষেত্র। সব সরকারি ডেটা একটি মাত্র টেবিলে থাকে না; উদ্দেশ্য ও দায়িত্ব অনুযায়ী পৃথক ব্যবস্থা থাকে।

\subsection{ব্যবহারের উদাহরণ}
জন্মনিবন্ধনে ব্যক্তির জন্মসংক্রান্ত তথ্য ও নিবন্ধন পরিচিতি রাখা যায়। ভূমি ব্যবস্থায় জমির পরিচয়, মালিকানা ও সংশ্লিষ্ট নথির তথ্য থাকে। শিক্ষা ব্যবস্থায় প্রতিষ্ঠান, শিক্ষার্থী ও পরীক্ষার ফল এবং স্বাস্থ্য ব্যবস্থায় সেবা ও রোগীসংক্রান্ত তথ্য সংরক্ষণ করা যায়। জনশুমারির সংগঠিত উপাত্ত থেকে এলাকা অনুযায়ী জনসংখ্যার চিত্র তৈরি করে পরিকল্পনায় ব্যবহার করা হয়।

\begin{flushleft}\small
\begin{tabularx}{\linewidth}{|p{31mm}|X|}\hline
\chSixHeader ক্ষেত্র & ডেটা ও সম্ভাব্য ব্যবহার\\\hline
নাগরিক সেবা & নিবন্ধন পরিচিতি দিয়ে সঠিক রেকর্ড খোঁজা\\\hline
ভূমি প্রশাসন & জমি ও নথির পরিচয় মিলিয়ে সেবা দেওয়া\\\hline
শিক্ষা & প্রতিষ্ঠান ও পরীক্ষার ফলের তথ্য পরিচালনা\\\hline
স্বাস্থ্য & সেবার চাহিদা ও প্রাপ্ত সেবার তথ্য বিশ্লেষণ\\\hline
পরিকল্পনা & অঞ্চলভিত্তিক পরিসংখ্যান ও সম্পদ বণ্টনের সহায়ক তথ্য\\\hline
\end{tabularx}
\end{flushleft}


\textbf{ভুল মিলের একটি ফল:} একই নামের দুই আবেদনকারীর মধ্যে কেবল নাম দেখে সেবা বা ফি-এর রেকর্ড মেলালে ভুল ব্যক্তির আবেদন নিষ্পত্তি হতে পারে। নির্ভরযোগ্য পরিচিতি, তথ্য যাচাই ও সংশোধনের সুযোগ তাই জরুরি। এক দপ্তরের তথ্য আরেক দপ্তরে পাওয়া উপকারী হলেও সব কর্মীকে সব নাগরিকের ব্যক্তিগত তথ্য দেখানোর অনুমতি দেওয়া প্রয়োজন হয় না।

\subsection{সুবিধা}
\begin{enumerate}
\item \textbf{অনুসন্ধান ও যাচাই:} পরিকল্পিত ডেটাবেজে তথ্য খোঁজা ও যাচাই সহজ হয়।
\item \textbf{পুনঃসংগ্রহ কমানো:} একই তথ্য বারবার সংগ্রহের প্রয়োজন কমতে পারে।
\item \textbf{সেবার অগ্রগতি:} সেবার অগ্রগতি অনুসরণ করা যায়।
\end{enumerate}
\subsection{পরিচালনার দায়িত্ব}
\begin{enumerate}
\item \textbf{পরিচয় যাচাই:} শুধু একই নাম দেখে দুই ব্যক্তির রেকর্ড এক করা যাবে না; উপযুক্ত পরিচিতি ও যাচাই প্রয়োজন।
\item \textbf{পরিবর্তনের ইতিহাস:} কোনো ভুল সংশোধনের ক্ষেত্রেও অনুমোদিত পদ্ধতি ও পরিবর্তনের ইতিহাস রাখা গুরুত্বপূর্ণ।
\item \textbf{ব্যক্তিগত তথ্যের সুরক্ষা:} নাগরিকের ব্যক্তিগত তথ্য সব ব্যবহারকারীর জন্য উন্মুক্ত হওয়া উচিত নয়।
\item \textbf{নির্ভরযোগ্যতা:} প্রয়োজন অনুযায়ী অনুমতি, সুরক্ষিত আদান-প্রদান, ব্যাকআপ ও নির্ভরযোগ্যতা বজায় রাখতে হয়।
\item \textbf{প্রকাশের উদ্দেশ্য:} জনসাধারণের জন্য প্রকাশযোগ্য পরিসংখ্যান এবং কোনো ব্যক্তির বিস্তারিত পরিচয় আলাদা উদ্দেশ্যের তথ্য।
\item \textbf{সীমিত অধিকার:} যে কর্মীর যতটুকু তথ্য প্রয়োজন তাকে ততটুকু অধিকার দেওয়া সঠিক পরিচালনার অংশ।
\end{enumerate}

\subsection{সেবা দেওয়ার তথ্যপ্রবাহ}
একটি সেবা কেন্দ্রে আবেদন গ্রহণ, পরিচয় যাচাই, প্রয়োজনীয় নথি পরীক্ষা, সিদ্ধান্ত এবং সেবার অগ্রগতি আলাদা তথ্য হিসেবে রাখা যায়। প্রতিটি আবেদনের ApplicationID থাকলে নাগরিকের নাম পুনরাবৃত্ত হলেও নির্দিষ্ট আবেদন খুঁজতে সুবিধা হয়। একই নাগরিকের একাধিক আবেদন থাকতে পারে; তাই CitizenID ও ApplicationID-এর কাজ আলাদা।

নাগরিকের ঠিকানা সংশোধনের সময় তথ্যটি কোন কোন সেবায় ব্যবহৃত হচ্ছে, তা বিবেচনা করতে হয়। সংশোধনের অনুমোদন, সময় ও দায়িত্বপ্রাপ্ত কর্মীর পরিচয় রাখলে পরে ভুলের কারণ খুঁজে পাওয়া সহজ হয়। প্রতিবেদনে মোট আবেদন ও নিষ্পত্তির হার দেখানো গেলেও সবার ব্যক্তিগত নথি প্রকাশের প্রয়োজন নেই। জনসেবার দ্রুততা, তথ্যের শুদ্ধতা ও ব্যক্তিগত তথ্যের সুরক্ষা একসঙ্গে নকশায় রাখতে হয়।

উদাহরণস্বরূপ, মাসে 500টি আবেদন পাওয়া এবং 450টি নিষ্পত্তি হওয়ার উপাত্ত থেকে 90\% নিষ্পত্তির হার পাওয়া যায়। কিন্তু বাকি 50টির কারণ জানতে অবস্থা ও প্রয়োজনীয় মন্তব্যের রেকর্ড লাগে। তাই শুধু মোট সংখ্যা নয়, উদ্দেশ্য অনুযায়ী প্রয়োজনীয় বিস্তারিত তথ্যও নকশায় রাখতে হয়।

\subsection{গুরুত্বপূর্ণ প্রশ্ন}
\textbf{অনুশীলন:} একই নামের দুই নাগরিকের রেকর্ড আলাদা রাখতে কী প্রয়োজন? জনশুমারির উপাত্ত থেকে একটি অর্থবহ তথ্যের উদাহরণ দাও। সরকারি ও কর্পোরেট ডেটাবেজের উদ্দেশ্য আলাদা হলেও কোন ব্যবস্থাপনা-নীতিগুলো উভয়ের জন্য প্রয়োজন?


